\begin{explanationbox}
	\textbf{\textcolor{CExplanation}{一句话直觉}}
	等腰三角形条件转化为两点距离相等，即中垂线或圆，再与动点轨迹联立。

	\medskip
	\textbf{\textcolor{CExplanation}{核心拆解}}
	\begin{description}[style=nextline, leftmargin=2.8em, labelsep=0.5em, itemsep=0.45em]
		\item[\textbf{要点一（等腰条件转化）：}]
			三角形 $ABP$ 中任两边相等，写出距离等式，如 $PA=PB$，$PA=AB$，$PB=AB$。
		\item[\textbf{要点二（三种分类）：}]
			不确定哪两边相等时，必须逐一考虑三种情形，不可遗漏。
		\item[\textbf{要点三（联立求解）：}]
			每个距离等式均为 $x,y$ 的方程（直线或圆），与轨迹 $f(x,y)=0$ 联立得候选点。
		\item[\textbf{要点四（排除共线）：}]
			若 $A,B,P$ 共线，则退化，需剔除。
	\end{description}

	\medskip
	\textbf{\textcolor{CExplanation}{几何本质（轨迹并集）}}
	满足 $PA=PB$ 的点 $P$ 轨迹是 $AB$ 的中垂线；满足 $PA=AB$ 的点 $P$ 轨迹是以 $A$ 为圆心、$|AB|$ 为半径的圆；满足 $PB=AB$ 的点 $P$ 轨迹是以 $B$ 为圆心、$|AB|$ 为半径的圆。三者的并集即为所有可能位置。

	\medskip
	\textbf{\textcolor{CExplanation}{代数意义（方程联立）}}
	将上述三个方程分别与 $f(x,y)=0$ 联立，转化为方程组求解，得到 $P$ 的坐标。每个方程均为二次以下，可代数求解。

	\medskip
	\textbf{\textcolor{CExplanation}{考点价值（存在性讨论）}}
	常用于圆锥曲线中探究是否存在点 $P$ 使三角形等腰，常涉及参数范围、交点个数等。常与弦长、垂直、斜率结合。

	\medskip
	\textbf{\textcolor{CExplanation}{顿悟点}}
	\textbf{分类要全，共线要舍。} 不要只考虑 $PA=PB$ 而忘记腰与底边相等的情形。

	\medskip
	\textbf{\textcolor{CExplanation}{使用场景}}
	\begin{itemize}[leftmargin=*, itemsep=0.5em, topsep=0.4em]
		\item \textbf{场景一（定点定轨迹）：} 已知 $A,B$ 和曲线 $C$，求 $C$ 上点 $P$ 使 $\triangle ABP$ 等腰。
		\item \textbf{场景二（参数问题）：} 已知 $A,B$ 和含参曲线，问是否存在参数使等腰三角形存在。
	\end{itemize}

\end{explanationbox}
